\begin{frame}[t]{对本项目的启示：身份映射、迁移与受控比较}
\bodyfont
\begin{columns}[T,onlytextwidth]
\begin{column}{.47\textwidth}
\subhead{可以借鉴的部分}
\begin{itemize}\setlength\itemsep{.14cm}
\item 可靠的基因--蛋白 ID 映射，加模态专用数值处理。
\item 从小规模处理矩阵开始，先验证非配对预训练收益。
\end{itemize}
\end{column}
\begin{column}{.49\textwidth}
\subhead{仍需本项目回答的问题}
\begin{itemize}\setlength\itemsep{.14cm}
\item 只给对照与药物时，RNA 能否改善蛋白变化预测？
\item 顺序迁移与非配对联合训练，哪种更能跨条件泛化？
\end{itemize}
\end{column}
\end{columns}
\vspace{.10cm}
\renewcommand{\arraystretch}{1.15}
\begin{tabularx}{\textwidth}{@{}P{3.5cm}Y@{}}
\toprule\textbf{最小实验}&\textbf{控制条件}\\\midrule
蛋白单任务 / RNA 预训练再微调 / RNA+蛋白联合 & 同一蛋白训练/测试集、输入信息与训练预算；对照不变、均值、线性模型作为基线\\
聚类或相似性加权 & 作为额外消融；只用训练数据拟合，不用测试蛋白响应寻找训练配对\\\bottomrule
\end{tabularx}
\source{本页是根据 [22--25] 形成的本项目研究建议，不是文献已证明的结论。分开报告“RNA 已见该药的跨模态迁移”与“两模态均未见药的外推”。}
\end{frame}
